\scanpage{91}
\mspage{44}

$\exists$ generator $g\in T$ (i.e. $g\in T^*$) and an order \inserted{$i_1$--$i_\nu$}
on the set $I$, \deleted{\illegible{2}} conjugates
$\unsure{l_i'|l_{i_\alpha}'}$ of the $\kappa_{i_\alpha}(g)$, and $g\in\mathbb N$ and $x_\alpha,y_\alpha\in\pi$
$(1\leq\alpha\leq g)$ such that we have
\[
\left(\prod_1^g[x_\alpha,y_\alpha]\right)\prod_1^{\nu}l_i=1
\]
and that this be a defining relation
of $\pi$.

There is \inserted{however} a grain of salt for $I=\varnothing$
(in which case $T$ serves no purpose \inserted{apparently} (in the
groupoid description of $(X,S)$ \unsure{$\simeq$|$=$} $(X,\varnothing)=X$)
one must then, \unsure{\deleted{in}} place of the
``Kummerian'' data \inserted{$\kappa$}, give oneself
\[
H^2(\pi,\mathbb Z)\overset{\kappa}{\simeq}T
\]
\margin{One must introduce \inserted{\unsure{therefore}} \unsure{this} as an additional datum,
in the definition of groups with loops the exclusion of the cases $X_0\simeq S^2$ (in the case of empty $S$) is \unsure{then} particularly \unsure{convincing}.}

Finally, one must \deleted{\illegible{1}}
(\unsure{even|$=$} with this grain of salt) exclude the case $X=S^2$, $S=\varnothing$
(as a connected component) --- i.e.
on the groups side, the case of a \inserted{component of $\Pi_U$ where}
$\pi=(1)$ \deleted{\illegible{2} $I=\varnothing$} and $I=\varnothing$.
[If I remember correctly, there is no reason
to exclude $(S^2,\mathrm{pt})$ (i.e. $X\setminus S=U\simeq\mathbb R^2\simeq\mathbb E^1_{\unsure{\mathbb C|e}}$)
\margin{90}

\scanpage{92}

where nevertheless \unsure{we have} $\pi=0$. But except in these
two cases, \deleted{the homomorphisms $\kappa_i:$} (corresponding
to the case $g=0$, $\nu=1$) the homomorphism $\kappa_i:T\longrightarrow\pi$
is injective. If $T=\mathbb Z$ the datum of the $\kappa_i$
is equivalent to that of elements $l_i\in\pi$, and we
recover the usual definition of
loop structures.

This is made explicit in the thesis of
Yves Ladegaillerie --- what is missing there, is
(among other things) the consideration of arrows
between $(X,S)$ other than homeomorphisms
(modulo isotopies)\,; for example maps \deleted{étale}
\deleted{\illegible{3}} $X'\xrightarrow{f}X$ such that $S'=f^{-1}(S)$
[and that $X'$ be étale over $X\setminus S$, if one insists\unsure{!}],
which reduces to the preceding case
--- the \inserted{more} general case where one assumes only
that $X'$ is a ramified covering
of $X$ (possibly ramified also
outside $S'$) and $S'\supset f^{-1}(S)$ (but $S'$
possibly larger) would call for
\margin{91}

\scanpage{93}
\mspage{45}

a careful study, with an
ad-hoc notion of isotopy ....

Another important direction (notably
for the study \unsure{of the|of the} unoriented,
nonorientable case) is the introduction of
finite groups of homeomorphisms, not
necessarily respecting the orientation.
To treat the case of a change of orientation,
note that in the groupoid description
$(\Pi_{D^*}\longrightarrow\Pi_U,\kappa)$ of a $(X,S)$ that is $T$-oriented,
the passage to the opposite orientation is expressed by
keeping $\Pi_{D^*}\longrightarrow\Pi_U$ as is, and
replacing $\kappa$ by $\bar\kappa=\kappa^{-1}$
\[
\bar\kappa(\unsure{\zeta|\xi})=\kappa(-\unsure{\zeta|\xi})=\kappa(\unsure{\zeta|\xi})^{-1}
\]
(If \inserted{$T=\mathbb Z$, on} the system $(\pi,(l_i))$, this amounts to
replacing the $l_i$ by the $l_i^{-1}$ ...)
\margin{and ditto (if $I=\varnothing$) for $\kappa:T\simeq H^2(\pi,\mathbb Z)$
\unsure{replaced} by $-\kappa$.}

Let then $\Gamma$ be a group (a priori
\unsure{not} necessarily finite) which operates on $(X,S)$, hence
on $U=X\setminus S$ \inserted{and on $\Pi_U$, $\Pi_{D^*}$.} We then find
mixed \inserted{fundamental} groupoids \deleted{\illegible{1}} by the well-known construction
\margin{92}

\scanpage{94}

\diagram{\begin{tikzcd}[ampersand replacement=\&]
\Pi_{D^*,\Gamma} \arrow[r] \& \Pi_{U,\Gamma} \arrow[r] \& \Pi_{e,\Gamma}
\end{tikzcd}}
category of $\Gamma$-torsors
i.e. of $1$-connected objects
in $\Gamma$-Sets

corresponding to the functors in the opposite direction

\diagram{\begin{tikzcd}[ampersand replacement=\&]
\text{\unsure{étale|étale} $\Gamma$-coverings of $D^*$} \& \text{\unsure{étale|étale} $\Gamma$-coverings of $U$} \arrow[l] \& \text{$\Gamma$-Sets} \arrow[l]
\end{tikzcd}}

The connected components of $\Pi_{D^*,\Gamma}$ correspond
to the \uline{orbits of $\Gamma$ in $S$} \deleted{\illegible{1}} $\simeq\pi_0(\Pi_{D^*})$,
likewise for those of $\Pi_{U,\Gamma}$. The \deleted{\unsure{topological}}
connected $\Pi_{U,\Gamma}$ \inserted{i.e. the connected topos $\widehat\Pi_{U,\Gamma}$,}
is the one \unsure{where|with} $\Gamma$ transitive
on $\pi_0(U)\simeq\pi_0(X)$ --- at the cost of replacing
$X$ by a connected component $X_0$, and $\Gamma$
by the \unsure{subgroup|subgroup} $\Sigma\subset\Gamma$ that stabilizes it,
we would be reduced (for the study of the
topoi $\widehat\Pi_{U,\Gamma}$ and of morphisms of topoi
\diagram{\begin{tikzcd}[ampersand replacement=\&]
\widehat\Pi_{D^*,\Gamma} \arrow[r,dash] \& \widehat\Pi_{U,\Gamma} \arrow[r] \& (\unsure{\widehat\Pi_{e,\Sigma}} \arrow[r] \& )\widehat\Pi_{e,\Gamma}
\end{tikzcd}}
to the case of $(X_0,S_0,\Sigma)$. But \unsure{the analogy|the analogue}
\margin{93}

\scanpage{95}
\mspage{46}

that I have in mind with the ``arithmetic'' case
makes this reduction inopportune in the
general case (for in the arithmetic \unsure{\deleted{case}},
we do not restrict ourselves either to the case where
$\Sigma=\Gamma$ i.e. $\mathcal E\longrightarrow\Gamma$ surjective, i.e. \inserted{the connected component}
$X$ \uline{geometric-}
\uline{ally} connected over $K$). Note \inserted{here} that
$(X,S,\Gamma)$ is recovered from the
$(X_0,S_0,\Gamma_0=\Sigma)$ as an amalgamated sum
\[
X=X_0\,\unsure{\wedge|\widehat\times}_{\Sigma}\Gamma\quad ....
\]

When we express (for $\Pi_{\unsure{U|X},\Gamma}$ connected)
the situation in terms of group theory,
we therefore find \deleted{\illegible{1}} a mixed
fundamental group \deleted{$\mathcal E_X$}
\[
\mathcal E_{\unsure{U|X},\Gamma}=\pi_1(U,\Gamma)
\]
and a homomorphism
\[
\mathcal E_{U,\Gamma}\longrightarrow\Gamma
\]
--- (Surjective iff \deleted{$\Gamma$ transitive on $\pi_0(U)=\unsure{\Gamma/\Sigma}$}
\inserted{$U$ connected (we have $\Gamma/\Sigma\simeq\pi_0(U)$)})
\inserted{finally a set of indices}
\deleted{corresponding to \illegible{1} $\pi_0(U)/\Gamma$}
$I$ ($\simeq\unsure{S}/\Gamma\simeq\pi_0(\Pi_{D^*,\Gamma})$) and
mixed fundamental groups
\margin{94}

\scanpage{96}

\[
\mathcal E_i\xrightarrow{\sim}\unsure{\pi_1|\Pi_1}(D_i^*,\Gamma)
\simeq\pi_1(D_{i_0}^*,\unsure{\Sigma_{i_0}|\Gamma_{i_0}})
\]
where $D_{i_0}^*$ is a connected component of the
punctured multidisk $D_i^*$ \inserted{(corresponding to the choice of an $\unsure{s_i}\in S$)}, and where $\unsure{\Sigma_{i_0}|\Gamma_{i_0}}\subset\Gamma$ is
its stabilizer (i.e. the stabilizer of
$s_i$ in $\Gamma$), which is \deleted{\unsure{a group}} (if $\Gamma_i$ \inserted{is finite and} operates
faithfully in a neighborhood of $s_i$) a cyclic
or dihedral group ... We therefore find, if $X$ is
$T$ \unsure{oriented|oriented}, an extension of $\Gamma_{i_0}$ by $T$,
the homomorphism $\mathcal E_i\longrightarrow\unsure{\Sigma_{i_0}|\Gamma_{i_0}}\subset\unsure{\Gamma|\Gamma_i}$ being induced
of course via $\mathcal E_{U,\Gamma}\longrightarrow\Gamma$ and $\mathcal E_i\longrightarrow\mathcal E_{U,\Gamma}$.

One must \inserted{also} link the action of $\Gamma$ on $X$ to
the orientation of $X$ --- for this we suppose
given a character
\[
\chi:\Gamma\longrightarrow\mathbb Z^*=\{\pm1\}
\]
and require that for $g\in\Gamma$, $g_X$ \unsure{preserve}
the orientation if $\chi(g)=1$, reverse it if $\chi(g)=-1$.
This implies that \deleted{\unsure{the}}\inserted{\unsure{we have an|the homomorphism or}} isomorphism
\[
\pi_1(\Pi_{D^*,\Gamma}/\Pi_{\unsure{e},\Gamma})
\overset{\kappa_\Gamma}{\xleftarrow{\sim}}T
\]
\margin{local system of the kernels of the
$\pi_1(\xi)\longrightarrow\pi_1(\varphi p(\xi))$
for $\xi\in\operatorname{Ob}\Pi_{D^*,\Gamma}$.}
\margin{95}

\scanpage{97}
\mspage{47}

$T$ being regarded as a local system
on $\Pi_{\Gamma,e}$ \deleted{\illegible{1}} i.e. as a $\Gamma$-group, thanks to the action
of $\Gamma$ via the character $\chi$. It amounts to the
same to say that \inserted{$\forall i\in I$} the map
\deleted{The datum}
\[
\kappa_i:T\xrightarrow{\sim}L_i=\operatorname{Ker}(\mathcal E_i\longrightarrow\Gamma)
\]
as a homomorphism from $T$ to $\mathcal E_i$ is
compatible with the action of $\mathcal E_i$ (operating
on $T$ via $\chi p_i$ ($p_i:\mathcal E_i\longrightarrow\Gamma$)), \illegible{1} and on
itself by inner automorphism\unsure{s} ...

If we exclude the case where $(X_0,\unsure{S_0})$
is $\simeq(S^2,\varnothing)$ or $(S^2,1\ \mathrm{pt})$ (i.e.
the case $\pi=0$), the $\kappa_i:T\longrightarrow\pi$ are injective,
hence also the $\mathcal E_i\longrightarrow\mathcal E_{U,\Gamma}$, hence the
$\mathcal E_i$ can be regarded as
\uline{subgroups} of $\mathcal E_{U,\Gamma}$.

\unsure{\S} Here it \deleted{\illegible{1}} would be particularly
inadvisable (\unsure{even|$=$} if we assume $\Sigma=\Gamma$ i.e.
$X=X_0$ i.e. $X$ connected) to suppose that
the $\Gamma_i$ \unsure{be|are} $=\Gamma$ i.e. that the $s\in S$ are fixed
by $\Gamma$ !
\margin{96}

\scanpage{98}

Since the center of $\pi$ is \deleted{\illegible{1}}reduced to $1$
(if we except \deleted{\unsure{resp.}} the case $(X_0,S_0)\simeq(S^2,\text{two points})$
i.e. $U_0\simeq\mathbb C^*$ \inserted{--- case of the annulus ---}), the datum of an extension
of $\Sigma$ ($\subset\Gamma$) by $\pi$ amounts (up to a unique \unsure{isomorphism})
to the datum of a homomorphism
\[
\unsure{\Sigma|\Gamma}\longrightarrow\mathrm{Aut\,ext}(\pi)\quad\deleted{\text{loop structure}}
\]
\deleted{from $\Sigma$ into the ``Teichmüller group''
relative to $\pi$ ($\simeq\pi_1(X_0)$)} --- $\mathcal E$ is reconstructed
as the inverse image of the extension
\diagram{\begin{tikzcd}[ampersand replacement=\&]
1 \arrow[r] \& \pi \arrow[r] \& \mathrm{Aut}(\pi) \arrow[r] \& \mathrm{Autext}\,\pi \arrow[r] \& 1
\end{tikzcd}}.

But here the outer \unsure{automorphisms} relative
to the $\alpha\in\Sigma$ \deleted{considered are} \uline{respect the loop structure}
\deleted{\illegible{1}} of $\pi$, modulo the
sign $\chi(\alpha)$ --- i.e. for $g\in\mathcal E$, and $s\in S$
$\exists\ \alpha\in\pi$ \inserted{and $s'\in S$ ($s'$ unique !)} such that
\[
\deleted{\unsure{\varphi|g}}\quad
\operatorname{int}(g)\kappa_s(\unsure{\zeta|\xi})
=\operatorname{int}(\alpha)\kappa_{s'}(\chi(p(g))\unsure{\zeta|\xi})
\]
$\forall\unsure{\zeta|\xi}\in\unsure{T|\mathbb Z}$, i.e.
\[
\operatorname{int}(\alpha^{-1}g)\kappa_s(\unsure{\zeta|\xi})
=\deleted{\unsure{\chi}}\unsure{\kappa_i|\kappa_{s'}}(\chi(p(g))\unsure{\zeta|\xi})
\]
\deleted{which implies that $\alpha^{-1}g$ normalizes $L_s=\kappa_s(T)$}
\margin{97}

\scanpage{99}
\mspage{48}

\deleted{(so probably $\alpha^{-1}g\in\mathcal E_s$ ?), and
operates on it \illegible{2}}
Thus, the operation of
$\Sigma$ on $S_0$ (hence of $\Gamma$ on $S=S_0\wedge_{\Sigma}\Gamma$)
is known, from the outer operation of $\Sigma\subset\Gamma$ on
$\pi$ endowed with its ``$T$-loop'' structure \deleted{\illegible{1}}
(i.e. the outer homomorphisms
$\kappa_s:T\longrightarrow\pi$) --- hence also $I=S_0/\Sigma\simeq S/\Gamma$.

Can we reconstruct $\mathcal E_i$ \inserted{($i\in I$)} from the
\deleted{\unsure{if} $i=s\bmod\Sigma$}
extension structure? We see, \deleted{that}
\deleted{$\mathcal E_{i_0}\subset\operatorname{Norm}_{\mathcal E}(L_s)$}
by virtue of the choices made,
that if $i\in I$ (hence $i$ an orbit of
$\Sigma$ in $S_0$) $\exists\ s\in i$ such that $\operatorname{Ker}(\mathcal E_i\longrightarrow\Gamma)$
is none other than $L_s=\kappa_s(T)$, and we
therefore have $\mathcal E_i\subset\operatorname{Norm}_{\mathcal E}(L_s)$, But
we have
\[
\operatorname{Norm}_{\mathcal E}(L_s)\cap\pi
=\operatorname{Norm}_{\pi}(L_s)=L_s
\]
$\unsure{=|\ne}\mathcal E_i\cap\pi$, and on the other hand the image of
$\operatorname{Norm}_{\mathcal E}(L_s)$ in $\Sigma\subset\Gamma$ is not \deleted{other than} \inserted{$\subset$}
$\Sigma_s=$ stabilizer \deleted{of} \inserted{$s$ in} $\Sigma$
$=$ Image of $\mathcal E_i$ in $\Sigma$,
\unsure{and|hence} in summary
\[
\mathcal E_i=\operatorname{Norm}_{\mathcal E}(L_s)
\]
\margin{98}

\scanpage{100}

Conversely, the datum of $\mathcal E\longrightarrow\Gamma$
and of the $\mathcal E_i$, $\kappa_i:T\longrightarrow\mathcal E_i$ gives back the
loop structure of $\pi$, by taking
the $\kappa_i$ and all their \inserted{distinct} outer
conjugates by the $\alpha\in\Sigma$ (modified by $\chi(\alpha)$ ...).
Thus the \inserted{datum of the} situation $\mathcal E\xrightarrow{p}\Gamma$, $\mathcal E_i\subset\mathcal E$
\deleted{equivalent} (family of subgroups, each
\unsure{defined} modulo conjugation \uline{in $\mathcal E$}) is equivalent
to the datum of

a) $\pi=\operatorname{Ker}p$, with its \inserted{$T$-}loop structure
(\deleted{finite family of \unsure{conjugacy classes}}\inserted{finite set of outer homomorphisms from}
$T$ to $\pi$)

b) A \unsure{sub-}\deleted{group} group $\Sigma\subset\Gamma$, and
a homomorphism
\[
\Sigma\longrightarrow\text{extended Teichmüller of }\pi
\]
(outer automorphisms of $\pi$,
respecting the loop structure
modulo sign
compatible with the character $\chi|\Sigma$ and
the ``sign'' character on extended Teichmüller.
In fact, $(X,S,\Gamma)$, or $(U,\Gamma)$, defines
$\pi$ only as an outer group \inserted{with $T$-loops} on which
$\Gamma$ operates in a way compatible with $\chi$.
\margin{99}
